%% kapa: exam=1 round=21 year=2010 period=1 form=A subject=civil_law subject_ko=민법 seq=3 number=3
%% source: Q-Net 감정평가사 기출문제 (cst00309) artlSeq=5024181 file=extracted/1008/5024181/2010년 제21회 감정평가사 1차시험 기출문제/1교시.hwp
%% conversion: kordoc 4.24.51 -> tools/gen1.py (auto); review=proofread-sonnet; flags=none
\begin{question}{3}
조건과 기한에 관한 설명으로 옳지 않은 것은? (다툼이 있으면 판례에 의함)
\begin{choices}
  \item 조건이 선량한 풍속 기타 사회질서에 위반한 것인 때에는 그 조건 뿐만 아니라 법률행위 전부가 무효로 된다.
  \item 기성조건이 해제조건이면 그 법률행위는 무효이다.
  \item 기한 있는 법률행위의 당사자는 기한의 도래가 미정인 동안 기한의 도래로 인하여 생길 상대방의 이익을 해하지 못한다.
  \item 기한은 소급효가 없으며 당사자의 특약에 의해서도 소급효가 인정되지 않는다.
  \item 기한이익상실의 특약은 특별한 사정이 없는 한 정지조건부 기한이익상실의 특약으로 추정한다.
\end{choices}
\end{question}
\answer{5}
